\section{引言}
我们把这次作业当作一次从源程序到可执行文件的追踪练习。只看最后的运行结果，很难知道宏、类型检查、控制流、符号和库函数分别在哪个阶段发生变化。因此，小题1保留了预处理文件、token、tree、CFG、RTL、汇编、目标文件和链接结果，再结合源码逐项说明它们的作用。

小题2以同一份 SysY 程序为基础，分别编写 LLVM IR 和 RISC-V 汇编。两种实现使用相同的函数划分和测试输入，运行结果用于检查数组访问、循环、分支、函数调用和运行时接口是否一致。

\section{实验环境与总体方法}
成员A的小题1实验在 WSL2 Ubuntu 22.04 环境中完成，主要使用 GCC 11.4.0，并在词法分析阶段使用 Clang 14 输出 token 序列，目标平台为 x86-64。成员B的小题1实验在 Ubuntu 22.04 Docker 容器中完成，主要使用 Clang/LLVM 14.0.0、GCC 11.4.0、Graphviz 2.43.0；流程实验另观察 AArch64 目标汇编与 ELF 文件。成员A的 LLVM IR 实验以 x86-64 Linux 为运行目标，在 Ubuntu 22.04 Docker 容器中完成。RISC-V 部分使用课程提供的交叉工具链与 QEMU，目标为 RV64，并链接课程提供的 SysY RISC-V 运行时库。

实验分两条线进行。第一条线使用 C/C++ 程序，保存预处理文件、token、GCC tree、CFG、RTL、汇编、目标文件和可执行文件，并比较不同编译参数带来的变化。第二条线以共享 SysY 程序为准，分别完成 LLVM IR 和 RISC-V 版本，再通过汇编、链接和多组输入检查两种实现的输出。
